\chapter{Discussion} \label{sec:Discussion}

The system was evaluated as a proof of concept against the requirements of the problem statement (\Cref{sec:ProblemStatement}), verified through controlled bench and climate-chamber testing rather than full offshore qualification. This chapter interprets the results of \Cref{sec:Results} in that light. It first assesses, sub-question by sub-question, what the work has established and what remains open. It then turns to the central tension running through the sensor development, between the sensitivity that makes the printed sensor usable at the small strains of interest and the reproducibility this same sensitivity costs, before examining how much confidence the test setup itself allows in the reported numbers. The chapter closes by considering what these results mean for turning the prototype into a deployable product, and which directions of future work would do the most to close the remaining gaps.

\section{Assessment Against the Research Question}
The first sub-question asked how 3D printed sensor characteristics could be optimized for accurate strain transfer and sensitivity. The five-stage print optimization of \Cref{Baseline Single-Sensor Response} answers this concretely for the chosen conductive sensing material, PLA, raising the relative resistance change at \SI{800}{\micro\strain} from \num{0.08} to \num{0.38}, with the line depth setting responsible for the largest single gain. The resulting gauge factor of approximately \num{480}, and up to \num{750} for one arm of the assembled bridge (\cref{sec:strain_sensing_performance}), is more than two orders of magnitude above the value of a metallic foil gauge~\cite{metal_gauge}, demonstrating that usable sensitivity at strains far below those previously reported for this sensor type~\cite{G.O.A.T} is achievable through print-process control alone, without a change of sensing principle.

The second sub-question, on integrating a strain measurement system within a $30\times\SI{30}{\milli\meter}$ metallic surface, is answered by the assembled module itself. The Wheatstone bridge pair (\SI{25}{\milli\meter} by \SI{6.25}{\milli\meter}, \cref{Baseline Single-Sensor Response}), the PCB (\SI{20.83}{\milli\meter} by \SI{25.02}{\milli\meter}, \cref{sec:final_pcb}), and the \textonehalf AA battery holder (\SI{34.4}{\milli\meter}~$\pm$~\SI{0.5}{\milli\meter} by \SI{17.5}{\milli\meter}~$\pm$~\SI{0.5}{\milli\meter}, \cref{sec:battery}) together fit within the dimensional constraints of the nut surface. The drift characterization on a steel-plate mounting surface (\Cref{Material & Mounting-Surface Selection (Drift)}) further confirms that the sensor can be printed directly onto a metal surface, though this was not tested under strain. The sensor length of \SI{25}{\milli\meter} does not yet use the full \SI{30}{\milli\meter} available on the nut face, and extending the serpentine to span the full width would be a straightforward way to improve sensitivity further. Conversely, the battery holder's \SI{34.4}{\milli\meter} dimension exceeds the \SI{30}{\milli\meter} nut face, so unless this overhang is acceptable in the final mounting, a smaller battery holder should be found.

The third sub-question concerned wireless communication suited to acquiring data across the many bolted joints of a turbine. The Zigbee mesh network described in \Cref{sec:zigbee_network} addresses the scale of the deployment directly, and the signal-strength results of \Cref{sec:signal_strength} show that the encapsulated module sustains a link at ranges relevant to a turbine interior, reaching a maximum of \SI{20}{\meter} against \SI{25}{\meter} for the bare module (\cref{tab:DistanceRange}). The silicone potting therefore costs only about \SI{3}{\decibel} of range, recoverable with a single transmit-power step. Whether this link margin is sufficient through an entire turbine structure, however, was not tested in the field and remains open.

The fourth sub-question asked how ultra-low-power electronics could support autonomous operation beyond ten years. The measured sleep- and run-mode currents translate, via the charge budget of \cref{tab:charge_10y}, into an estimated ceiling of approximately 55 reports per day for a full decade within the derated capacity of the selected battery at the highest measured current draw (\SI{50}{\celsius}), with substantially more headroom at lower reporting rates or lower temperatures. This is a budget computed from measured currents and extrapolated over ten years, however, not a demonstrated ten-year lifetime. 

The fifth sub-question, on performance across operating temperature and long-term use, was characterized extensively. The individual sensor and the assembled bridge were both cycled over the full \SIrange{-20}{50}{\celsius} range (\Cref{Temperature Behaviour}, \cref{fig:FS_ADCogVout}), and the bridge was run continuously for approximately \SI{60}{\hour} under constant conditions (\cref{fig:Slope1}) and under a smaller temperature excursion of \SIrange{20}{25}{\celsius} (\cref{fig:sidsteTest}). These tests establish that the sensor and bridge remain functional and track applied strain throughout, but it was not possible to vary strain across the full \SIrange{-20}{50}{\celsius} operating range, so the sensor's response to temperature and to strain were characterized simultaneously only over the smaller \SIrange{20}{25}{\celsius} excursion of \cref{fig:sidsteTest}. Combined with the fact that, as the following sections of this chapter discuss, several aspects of the test setup limit how precisely the underlying temperature and long-term drift of the sensor itself can be separated from confounding effects, this sub-question is only partially closed.

The final sub-question asked how encapsulation could resist water, oil, and grease while preserving strain transfer. A two-part silicone compound was selected and shown to preserve the sensor's strain response after potting (\Cref{sec:encapsulation}) while attenuating the wireless link only slightly (\Cref{sec:signal_strength}). The encapsulation is implemented and shown not to compromise sensing, while only slightly compromising communication, but its resistance to water, oil, and grease ingress, and its durability over years of environmental cycling, was not tested directly, so this sub-question is answered at the level of feasibility rather than qualification.

Taken together, the sub-questions show that the core sensing and system-integration challenges posed by the research question have been demonstrated on the bench, while the questions of field performance, absolute long-term accuracy, and multi-year durability remain open, and are treated in turn in the remainder of this chapter.

\section{The Sensitivity–Reproducibility Trade-off} \label{sec:tradeoff}
The high sensitivity established in the preceding section and the reproducibility limitations discussed below share a common origin, and making it explicit accounts for the limitations of the sensor as much as for its central achievement. That sensitivity follows directly from operating the serpentine at the near-separation resting point of \Cref{Sensor Structure & the Small-Strain Operating Point}, where the asperity contacts between adjacent turns sit just short of losing contact, and this is the regime in which a strain of a few hundred microstrain still produces a resistance change large enough to measure. The same proximity to separation that supplies this sensitivity also leaves the baseline resistance of every printed element acutely sensitive to the print process, to temperature, and to time, which are precisely the influences the sensor must tolerate in service. This behavior is therefore best understood as a trade-off inherent to operating a constriction-resistance sensor at this point, rather than as a deficiency of the fabrication that a better printer would simply remove.

Three consequences follow. The first concerns the optimization itself, where the print-to-print variation in the asperity geometry is larger at this operating point than it would be for a sensor resting far from separation. The line depth, which fixes how far the nozzle is lowered toward the substrate and is the dominant setting in the five-stage sequence of \cref{fig:opt_pla}, is among the hardest to hold constant between prints, since the nozzle-to-substrate spacing varies from one print to the next by amounts on the order of micrometers, and at the near-separation operating point a variation even this small shifts the resting contact enough to move the sensing characteristics noticeably. The practical difficulty is therefore not seeing which settings help, but reproducing the exact sensing characteristics from one print to the next, so that while the direction of each optimization is evident, the precise response of a given element is harder to guarantee. A study under tighter process control would likely both sharpen the smaller effects and narrow this run-to-run spread.

The second concerns the temperature compensation and the confidence that can be placed in it. The Wheatstone half-bridge of \Cref{Wheatstone Half-Bridge Design} rejects the shared temperature response of its arms only to the extent that those arms are matched, and while two elements can in principle be printed to the same characteristics, the print-to-print variation present in this work makes such close matching difficult to achieve in practice. A bridge assembled from mismatched arms balances about its own offset rather than about zero, and a residual temperature sensitivity then survives compensation in proportion to that mismatch, as \cref{fig:sidsteTest} shows directly, where the output envelope of the deliberately mismatched bridge tracks the ambient temperature rather than remaining flat. Improving the matching of the arms is therefore one of the clearest routes to a better-compensated bridge, and the difficulty is a recognized one for 3D printed Wheatstone bridges in general~\cite{3DprintWheatstoneStrain} rather than particular to this sensor. The mismatch is not the whole story, however, since the fixture-induced strain discussed in \Cref{sec:reliability} couples into the bridge output during a thermal sweep and cannot be separated from any temperature dependence of the sensor itself, so the imperfect compensation observed here cannot be attributed to the arm matching alone.

The third concerns the path to a product. As the process stands, each bridge carries its own offset and its own residual temperature coefficient, so a single design cannot be calibrated once and copied identically across units, and per-unit calibration or field commissioning becomes the more realistic manufacturing model in place of a single type-calibration applied uniformly. Per-unit calibration and the binning of sensors by measured characteristics are in any case established practice elsewhere in the sensor industry, so nothing in the printed sensor precludes either approach. This need is a consequence of the current print-to-print variation rather than a permanent feature, however, and if the printing were made reproducible enough for a sensor to come out the same way each time, a single type-calibration applied once to the design could replace the per-unit step.

These consequences are also bound to the equipment available for the project. The printer used is capable, but the achievable control over extrusion rate and the nozzle-to-substrate spacing is nonetheless limited by it, and a more capable machine, or one purpose-built for this deposition, would be expected to narrow the asperity-level variation at its source, improving both the reproducibility of the optimization and the matching of the bridge arms. This avenue is taken up in \Cref{sec:improved_test} rather than pursued here.

\section{Reliability and Validity of the Measurements}
The characterization above rests on a test bench built for this project rather than on calibrated commercial test equipment, and several aspects of that bench limit how far its results can be trusted as absolute measurements of sensor performance.

The actuation is open loop, and the applied elongation is computed from the number of microsteps commanded to the stepper driver (\cref{eq:microstep}), not measured independently by a displacement or force sensor at the point of contact with the device under test. Backlash in the lead screw, friction in the linear stage, and compliance in the fixturing are therefore unverified, so the nominal strain values reported throughout this thesis should be read as the commanded strain rather than as a strain established by an independent, traceable reference.

Repeated mounting and dismounting of a sensor, and the break-in and contact-settling behavior described in \Cref{Sensor Structure & the Small-Strain Operating Point}, further mean that even nominally the same physical sensor is not perfectly reproducible run to run, an effect visible in the continued cycle-to-cycle drift of \cref{fig:Slope2} despite constant temperature and humidity over the run.

The test rig itself is not thermally inert. Over a full temperature sweep the lead screw, frame, and fixturing of the bench expand and contract along with the device under test, and because the sensor is, by design, sensitive enough to register a few hundred microstrain, this rig-induced strain couples directly into the measured signal. This is the most likely explanation for the large, non-repeating excursions of the bridge output seen across the thermal cycle of \cref{fig:FS_ADCogVout}, where the output at a given temperature differs by tens of millivolts between successive visits to the same setpoint, a pattern more consistent with mechanical strain from the expanding fixture than with a fixed residual temperature dependence of the bridge. Because the sensor was not mechanically decoupled from this fixture-induced strain, the test can neither confirm nor rule out a temperature dependence of the bridge intrinsic to the sensor itself, limiting the strength of the conclusions that can be drawn from it, as already noted in \Cref{sec:strain_sensing_performance}.

A related limitation follows from the fact that strain could not be applied inside the climate chamber across the full range of temperature extremes explored. The effect of exposure to temperatures from \SI{-40}{\celsius} to \SI{80}{\celsius} on the sensor was therefore established by first exposing the sensor to a given temperature and only afterward straining it at room temperature (\cref{fig:PLA_after_temp_limits}), rather than by straining it while held at that temperature. This shows that exposure alone does not degrade the sensor's subsequent response, but it does not demonstrate that the sensor performs identically while simultaneously strained and held at a temperature extreme, a gap that matters most at the coldest and hottest ends of the intended \SIrange{-20}{50}{\celsius} operating range.

None of these limitations were addressed within the scope of this thesis, since doing so would have required a calibrated load frame with independent displacement and force readout, and a fixture capable of applying strain while thermally isolated from the surrounding environment, both beyond the equipment available for a project of this timeframe. With such equipment, discussed further in \Cref{sec:improved_test}, the absolute accuracy of the reported strain values, and the separation of genuine sensor drift from rig and fixture effects, could be established with considerably more confidence.

\section{Implications for a Deployable Product}
Read together, the results of this chapter place the system as a demonstrated proof of concept rather than a qualified product, and it is worth being explicit about where that boundary currently lies.

The ten-year operating life central to the research question is, at this stage, a budget rather than a demonstrated lifetime. \Cref{tab:charge_10y} shows that the measured currents are compatible with a decade of operation at realistic reporting rates, but no module has been run, or can practically be run within a master's thesis, for anything close to ten years. Likewise, the drift and temperature characterization of \Cref{sec:strain_sensing_performance} spans hours to a few days, not years, and the encapsulation has been shown to preserve sensing and RF performance immediately after potting (\Cref{sec:encapsulation}, \Cref{sec:signal_strength}) but has not been exposed to prolonged environmental cycling or ingress testing. What ultimately determines whether the system can detect a real loss of bolt preload in the field is the absolute long-term strain accuracy and drift of the sensor over months to years, which is precisely the quantity the present data cannot yet bound, for the reasons discussed above.

Against this, the results also support the core premise that motivated the design in the first place. The unit cost of approximately \SI{80}{DKK} (\cref{tab:cost_summary}) is low enough, relative to the value of monitoring a single high-consequence bolted joint, to make a one-module-per-bolt deployment plausible at the scale of a modern turbine's several thousand bolts~\cite{6000bolts}, in clear contrast to the proprietary, cost-constrained monitoring hardware discussed in \Cref{sec:Introduction}. Combined with the demonstrated sensitivity, wireless architecture, and low-power operation, this suggests that the remaining gaps are ones of qualification and field validation rather than of the underlying concept.


\section{Future Work}
The work presented in this thesis advanced a compact wireless strain-sensing system from concept to a verified prototype, but several avenues remain open before it could be deployed as a finished product. This section outlines the main directions for further development of the system, covering the sensor, the hardware, the wireless communication, the firmware, and the long-term qualification of the design.
 
\subsection{Full Bridge and Reference Arm Placement}
The present sensor uses the Wheatstone half-bridge of \cref{Wheatstone Half-Bridge Design}, in which two active arms are mounted on the strained surface and two passive reference arms are held unstrained for temperature compensation. As shown in \cref{sec:Wheatstone_theory}, replacing the two passive arms with active elements responding in the opposite direction forms a full bridge, which ideally doubles the output swing once more and linearises the output in the relative resistance change. Moving to a full bridge would therefore roughly double the sensitivity and improve the rejection of common-mode and temperature effects, at the cost of two additional active elements and the more demanding matching this requires.

A practical way to realize this is to mount all four arms on the strained surface but orient the two reference arms at \SI{90}{\degree} to the strain axis. Aligned transversely, these arms respond primarily to the Poisson strain, which is of opposite sign and smaller magnitude than the axial strain seen by the active arms, so they contribute a partial active response in the opposing direction while remaining far less sensitive than the axial arms. The advantage of this arrangement is mainly thermal. With all four arms printed together and bonded to the same surface, they are held at the same temperature and experience the same thermal environment, which directly addresses the matching and equal-temperature assumptions that \cref{Wheatstone Half-Bridge Design} identifies as the limiting factors in the compensation. 

\subsection{Print-Process Control and Sensor Reproducibility}\label{sec:process_control}
The trade-off examined in \Cref{sec:tradeoff} traced the print-to-print variation of the sensor to small differences in the print, most directly the nozzle-to-substrate spacing, whose effect is amplified at the near-separation operating point at which the sensor is most sensitive. The sensors in this work were produced on the printer available for the project, whose control over the extrusion rate and this spacing sets a floor on how closely two elements can be made to match. A more capable machine, or one purpose-built for this deposition, with closed-loop extrusion control and a temperature-regulated build environment, would be expected to narrow this variation at its source, improving both the reproducibility of the print optimization and the matching of the bridge arms, and with it the common-mode rejection on which the temperature compensation depends.
 
The reproducibility of the print also determines how the sensor must be calibrated in production. As the process stands, each bridge carries its own offset and residual temperature coefficient, so a per-unit calibration, or a binning of printed elements by measured characteristics, is the realistic model rather than a single calibration applied to every unit. Should the printing be made reproducible enough for elements to emerge with closely matched characteristics, this per-unit step could be reduced to a single type-calibration of the design, which would simplify manufacture considerably.
 
\subsection{Sensor Geometry and Mechanical Repeatability}
Two further refinements concern the printed element itself. The serpentine currently spans \SI{25}{\milli\meter} of the \SI{30}{\milli\meter} available on the nut face, so extending it to the full width offers a direct route to additional sensitivity without altering the sensing principle. Separately, the hysteresis under repeated loading and the initial settling observed in the rest-stability measurements were attributed to viscoelastic relaxation of the substrate and to the gradual settling of the asperity contacts, so a stiffer or less viscoelastic substrate, or a mechanical break-in procedure applied before calibration, would be worth investigating as a means of improving the short-term repeatability. How these properties evolve over years of service is left to the accelerated life testing described in \Cref{sec:improved_test}.

\subsection{Field Testing and Wireless Communication}
The module was validated in the laboratory and climate chamber, but not in its intended environment. A natural next step is to deploy it on an operational wind turbine, where the metallic nacelle and tower, the rotating machinery, and the long internal signal paths create propagation conditions that cannot be reproduced on the bench. Such a deployment would establish the true range and link margin and reveal environmental effects, such as vibration and electromagnetic interference, that the laboratory testing could not capture.
 
Should this testing show that the \SI{2.4}{\giga\hertz} link does not provide sufficient range or margin through the turbine structure, a study of sub-\si{\giga\hertz} implementations would be worthwhile. A sub-\si{\giga\hertz} signal is less attenuated by the metallic structure and less affected by interference from WiFi and Bluetooth, which yields additional link margin for the same transmit power~\cite{silabs_subghz_priorities}. This margin can be spent in either of two ways. It can extend the range, or it can be traded back by lowering the transmit power while still closing the link. As the radio transmission is one of the largest contributors to the energy consumed in each measurement cycle, a lower transmit power reduces the energy per transmission and thereby extends the battery life. Such a change need not mean abandoning Zigbee, as it is one of the few protocols available in both \SI{2.4}{\giga\hertz} and sub-\si{\giga\hertz} versions. Alternatively, a dedicated sub-\si{\giga\hertz} protocol such as Z-Wave could be considered, at the cost of a smaller and more expensive device ecosystem~\cite{homey_zigbee_vs_zwave}. 

\subsection{Sensor Commissioning}
In the current implementation, each module is identified manually, by assigning it a name in Home Assistant as it joins the coordinator. This step is workable for a few modules, but it does not scale well to a larger deployment and is prone to human error, as the installer must correctly match each joining device to its intended name and physical location. A useful extension would be to streamline this commissioning step, for example by labeling each module with a barcode or QR code encoding its identifier, so that it can be scanned and registered to a known position during installation rather than named by hand. This would make the mapping between modules and locations faster, less error-prone, and better suited to deployments involving many modules.

\subsection{Temperature Compensation}
The characterization of the on-die temperature sensor presented in \cref{sec:temp_sensor} established that it provides sufficient precision to serve as a reference for compensating the residual temperature sensitivity introduced by imbalance in the Wheatstone bridge. A natural extension of this work is therefore the derivation and integration of a compensation model that expresses the measured bridge output as a function of die temperature, allowing the residual thermal drift to be removed in firmware without modification to the existing hardware. 

\subsection{Battery Monitoring}
The current design does not measure its own battery voltage, so the state of the battery is not available in the user interface. For a system intended to operate unattended for years, this is a significant limitation, as there is no way to detect a weakening battery before it fails. The measurements already show a strong temperature dependence of the battery voltage  \cref{fig:FS_bat}), which makes such monitoring especially valuable, since a module may approach its voltage limit at low temperature while still appearing healthy under normal conditions. A useful extension is therefore to add battery-voltage sensing to the module, allowing the battery state of each module to be reported alongside its measurements and surfaced in the interface.

\subsection{Encapsulation Validation}
The present encapsulation encloses the PCB and sensor unit in a silicone compound (\cref{sec:encapsulation}). This was found to attenuate the RF link only slightly (\cref{sec:signal_strength}) and performed well overall, but its protective properties were not formally validated. Future work should therefore characterize the encapsulation more thoroughly, including extended environmental testing, mechanical drop testing, and verification of resistance to oil, grease, and water ingress. In parallel, alternative encapsulants could be investigated, with the aim of finding a solution that is both more economical and less attenuating to the RF signal than the current silicone compound.

\subsection{Improved Test Methodology and Long-term Qualification} \label{sec:improved_test}
The characterization in this work was performed with the available laboratory equipment, which limited the precision of the applied mechanical load and the simultaneous control of strain and temperature. A more rigorous characterization would use a calibrated load frame, such as an Instron, which compresses a fixture around the bolt to place it under a controlled and repeatable tensile preload, which is then released, while strain and temperature are measured simultaneously. This would allow the sensor response, during both loading and relaxation, to be related directly to a known applied load across the operating range. Beyond the sensor itself, the durability of the system has not been established. Accelerated life testing and ingress testing of the encapsulation would quantify how the protection and the sensor response evolve over time and under repeated environmental cycling. A field trial in the intended environment would finally confirm that the system performs as expected under real operating conditions over an extended period. Together, these steps would move the system from a validated prototype towards a qualified, deployable product.